\section{Recognizing an affine presentation}\label{recognizing-an-affine-presentation}

We write an integral affine transformation as \((M,t)\), acting by
\(v\mapsto Mv+t\), with multiplication

\[ (M,t)(N,u)=(MN,t+Mu). \tag{1} \]

Let \(E\) be a group generated by finitely many such transformations, and
let \(Q\leq\operatorname{GL}_d(\mathbb Z)\) be its finite linear image.
Its translation kernel \(L\) is a subgroup of \(\mathbb Z^d\), invariant under
\(Q\). A translation is written \(t_v=(I,v)\).

\textbf{Lemma AL.1 (fixed Schreier lattice).} Choose a word lift \(\widehat q\in E\)
for each \(q\in Q\), with \(\widehat1=1\). Then \(L\) is generated by the
translations

\[\widehat{\overline s q}^{-1}s\widehat q,\qquad
s\text{ a distinguished generator},\quad q\in Q.\tag{2}\]

\emph{Proof.} The products have trivial linear image. Write a partial word
as \(\widehat q n\), with \(n\in L\). Left multiplication by \(s\) changes
this to \(\widehat{\overline s q}t n\), where \(t\) is the transition in
(2). Induction expresses a word of trivial linear image as a product
of these transitions; its final representative is \(\widehat1=1\).
Inverses introduce no extra generator: the transition for $s^{-1}$ at $q$ is the inverse of the transition for $s$ at $\overline{s}^{-1}q$. Thus the induction applies to words with positive and negative letters. Every transition itself is in \(L\), so they generate exactly \(L\).
\(\square\)

The normalization in (2) matters. The raw differences
\(s\widehat q\widehat{\overline s q}^{-1}\) need not generate \(L\)
without conjugation. Each is a conjugate of the corresponding transition
by \(\widehat{\overline s q}\). Consequently their \(Q\)-orbits generate
\(L\). In particular, if a proposed \(Q\)-invariant lattice contains every
raw difference, it contains every normalized transition and hence \(L\).
If exhibited raw differences also span the proposed lattice, equality
follows. This invariant-lattice version justifies the raw-difference
format of the unchanged earlier \(\{6,6\}\) checker {[}C66{]}. The new
\(\{6,4\}\) checker uses (2) directly.

This is a finite proof of a lattice identity. Containment in a proposed
lattice is checked on the normalized transitions, or on the raw differences
with the stated invariance condition. Equality follows when independent
integer combinations give a basis of that lattice, or when the greatest
common divisor of the full-rank minors gives its exact index. Sampling
finite quotients is not a replacement for this step.

\textbf{Lemma AL.2 (presentation bridge).} Let \(U\) be a finitely presented marked
group whose relators hold on the affine generators of \(E\). Suppose their
linear images define the same epimorphism to \(Q\). If an exact
Reidemeister--Schreier presentation proves
\(\ker(U\to Q)\cong\mathbb Z^d\), and the affine translation kernel \(L\)
has rank \(d\), then the marked epimorphism \(U\to E\) is an isomorphism.

\emph{Proof.} The restriction maps \(\ker(U\to Q)\) onto \(L\): an affine
translation has a preimage in \(U\), and its linear image is trivial.
This is a surjection between free abelian groups of rank \(d\). Tensoring
with \(\mathbb Q\) shows that its kernel has rank zero. A subgroup of a
free abelian group is torsion-free, so that kernel is trivial. The
remaining kernel of \(U\to E\) lies in this subgroup, and is therefore
trivial as well. \(\square\)

The hypothesis concerns the actual kernel presentation. Computing only
its abelianization would not establish it: a nonabelian group can have
free abelian abelianization of rank \(d\). Our applications retain the
exact rewriting and Tietze algorithms used to verify this finite
presentation step. It is a computer-assisted input, not a Lean proof.
